\documentclass[10pt,a4paper]{amsart}
\usepackage{amsmath,amssymb,amsthm,graphicx}
% Retained minimal preamble for the Stacks Project chapter more-algebra.tex (More on Algebra).
% The chapter file pinned-original-source.tex carries only \input{preamble} on its first line, so no
% preamble text lives inside the chapter itself.  The Stacks Project uses one single shared preamble for
% every chapter; the snapshot stages it identically as sources/stacks-chow/preamble.tex and
% sources/stacks-etale-cohomology/preamble.tex.  The block below is copied mechanically, line for line,
% from sources/stacks-chow/preamble.tex lines 14-18 (the Xy-pic setup) and 154-217 (the native theorem
% declarations together with the mathematical macros).  Nothing was added, removed or redefined; the
% wrapper class and the amsmath/amssymb/amsthm/graphicx base packages are supplied by the assembler.
% Chapter: more-algebra.tex -- More on Algebra.
% For commutative diagrams we use Xy-pic
\usepackage[all]{xy}

% We use 2cell for 2-commutative diagrams.
\xyoption{2cell}

% Theorem environments.
%
\theoremstyle{plain}
\newtheorem{theorem}[subsection]{Theorem}
\newtheorem{proposition}[subsection]{Proposition}
\newtheorem{lemma}[subsection]{Lemma}

\theoremstyle{definition}
\newtheorem{definition}[subsection]{Definition}
\newtheorem{example}[subsection]{Example}
\newtheorem{exercise}[subsection]{Exercise}
\newtheorem{situation}[subsection]{Situation}

\theoremstyle{remark}
\newtheorem{remark}[subsection]{Remark}
\newtheorem{remarks}[subsection]{Remarks}

\numberwithin{equation}{subsection}

% Macros
%
\def\lim{\mathop{\mathrm{lim}}\nolimits}
\def\colim{\mathop{\mathrm{colim}}\nolimits}
\def\Spec{\mathop{\mathrm{Spec}}}
\def\Hom{\mathop{\mathrm{Hom}}\nolimits}
\def\Ext{\mathop{\mathrm{Ext}}\nolimits}
\def\SheafHom{\mathop{\mathcal{H}\!\mathit{om}}\nolimits}
\def\SheafExt{\mathop{\mathcal{E}\!\mathit{xt}}\nolimits}
\def\Sch{\mathit{Sch}}
\def\Mor{\mathop{\mathrm{Mor}}\nolimits}
\def\Ob{\mathop{\mathrm{Ob}}\nolimits}
\def\Sh{\mathop{\mathit{Sh}}\nolimits}
\def\NL{\mathop{N\!L}\nolimits}
\def\CH{\mathop{\mathrm{CH}}\nolimits}
\def\proetale{{pro\text{-}\acute{e}tale}}
\def\etale{{\acute{e}tale}}
\def\QCoh{\mathit{QCoh}}
\def\Ker{\mathop{\mathrm{Ker}}}
\def\Im{\mathop{\mathrm{Im}}}
\def\Coker{\mathop{\mathrm{Coker}}}
\def\Coim{\mathop{\mathrm{Coim}}}

% Boxtimes
%
\DeclareMathSymbol{\boxtimes}{\mathbin}{AMSa}{"02}

%
% Macros for moduli stacks/spaces
%
\def\QCohstack{\mathcal{QC}\!\mathit{oh}}
\def\Cohstack{\mathcal{C}\!\mathit{oh}}
\def\Spacesstack{\mathcal{S}\!\mathit{paces}}
\def\Quotfunctor{\mathrm{Quot}}
\def\Hilbfunctor{\mathrm{Hilb}}
\def\Curvesstack{\mathcal{C}\!\mathit{urves}}
\def\Polarizedstack{\mathcal{P}\!\mathit{olarized}}
\def\Complexesstack{\mathcal{C}\!\mathit{omplexes}}
% \Pic is the operator that assigns to X its picard group, usage \Pic(X)
% \Picardstack_{X/B} denotes the Picard stack of X over B
% \Picardfunctor_{X/B} denotes the Picard functor of X over B
\def\Pic{\mathop{\mathrm{Pic}}\nolimits}
\def\Picardstack{\mathcal{P}\!\mathit{ic}}
\def\Picardfunctor{\mathrm{Pic}}
\def\Deformationcategory{\mathcal{D}\!\mathit{ef}}

% --- appended 2026-10-12 (Chief): the project own definitions of the reference/slogan/history
% display environments, copied mechanically from the shared Stacks preamble lines 8-12
% (\usepackage{verbatim} plus the three \newenvironment lines, which the chapter preamble keeps
% verbatim from the project preamble). Nothing added, removed or redefined: these environments
% carry no mathematical content and the project itself typesets them as comments. Without these
% three lines a counted statement whose own text uses them fails with Environment slogan undefined.
\usepackage{verbatim}
\newenvironment{reference}{\comment}{\endcomment}
\newenvironment{slogan}{\comment}{\endcomment}
\newenvironment{history}{\comment}{\endcomment}
\usepackage{cleveref}
\crefname{originallemma}{original lemma}{original lemmas}
\Crefname{originallemma}{Original lemma}{Original lemmas}
\makeatletter
\newlabel{lemma-K-injective-flat}{{13926}{source}{}{sourcepointer-1}{}}
\newlabel{lemma-K-injective-flat@cref}{{[originallemma][0][]13926}{[1][0][]source}{}{}{}}
\makeatother
\crefname{originallemma}{original lemma}{original lemmas}
\Crefname{originallemma}{Original lemma}{Original lemmas}
\makeatletter
\newlabel{lemma-hom-from-Af}{{25850}{source}{}{sourcepointer-2}{}}
\newlabel{lemma-hom-from-Af@cref}{{[originallemma][0][]25850}{[1][0][]source}{}{}{}}
\makeatother
\crefname{originallemma}{original lemma}{original lemmas}
\Crefname{originallemma}{Original lemma}{Original lemmas}
\makeatletter
\newlabel{lemma-map-identifies-koszul-and-cech-complexes}{{24363}{source}{}{sourcepointer-3}{}}
\newlabel{lemma-map-identifies-koszul-and-cech-complexes@cref}{{[originallemma][0][]24363}{[1][0][]source}{}{}{}}
\makeatother
\title{Along a flat ring map inducing an isomorphism modulo a finitely generated ideal, restriction is an equivalence of derived complete derived categories}
\author{The Stacks Project Authors}
\date{stacks-project:more-algebra.tex:local-snapshot-20261012; GNU FDL 1.2 or later}
\begin{document}
\maketitle
\setcounter{subsection}{1}
\noindent\textit{Editorial note.} All delivered mathematical text is the author's own, copied from the pinned Stacks Project chapter file more-algebra.tex; no statement or proof step was rewritten, reordered or removed.\par\medskip
\noindent\textit{Editorial note.} The chapter file pinned-original-source.tex carries only \textbackslash{}input\{preamble\} on its first line, so no preamble text lives inside the chapter. The retained file sources/stacks-more-algebra/retained-preamble.tex is a mechanical, line-for-line copy of the project's single shared preamble as staged in the snapshot; nothing was added, removed or redefined and no package was added.\par\medskip
\noindent\textit{Editorial note.} The source URL is the source-level Stacks Project URL, as required when more than one card is produced from one source; the per-claim tag is carried by claim\_key.\par\medskip
\noindent\textit{Editorial note.} The author writes 'Some details omitted.' at proof line 26662 and leaves the final sentence of the essential-surjectivity paragraph without a full stop at proof line 26669; both are reproduced unchanged. The three references inside the counted proof span are recorded as non-load-bearing navigation pointers.\par\medskip
\noindent\textit{Allowed standard prerequisites (editorial).}\begin{itemize}
\item Derived complete objects D\_comp(A, I) and the vanishing of Ext out of the localizations A\_f for f in I. the hypothesis at statement lines 26636-26637 and used at proof lines 26648-26650. Is what makes the Cech resolutions of A and B quasi-isomorphisms against the given K-injective representative.
\item Extended alternating Cech complexes and the quasi-isomorphism comparing them for A and B. invoked at proof lines 26644-26646. Is the tool that transfers information from A to B along the flat map.
\item K-injectivity, including its preservation under the restriction of scalars. used at proof lines 26664 and 26676-26677. Lets Hom\_A(B, -) compute the correct derived objects and lets a B-module complex serve as an A-module complex.
\end{itemize}\medskip
\ifdefined\hypertarget\hypertarget{sourcepointer-1}{}\fi
\noindent\textit{Source reference.} Original lemma 13926: sources/stacks-more-algebra/pinned-original-source.tex, source lines 13926--13926. Same-chapter navigation pointer; the statement it names is standard for the section and the argument that consumes it stays inside the counted proof span.\par\medskip
\ifdefined\hypertarget\hypertarget{sourcepointer-2}{}\fi
\noindent\textit{Source reference.} Original lemma 25850: sources/stacks-more-algebra/pinned-original-source.tex, source lines 25850--25850. Same-chapter navigation pointer; the statement it names is standard for the section and the argument that consumes it stays inside the counted proof span.\par\medskip
\ifdefined\hypertarget\hypertarget{sourcepointer-3}{}\fi
\noindent\textit{Source reference.} Original lemma 24363: sources/stacks-more-algebra/pinned-original-source.tex, source lines 24363--24363. Same-chapter navigation pointer; the statement it names is standard for the section and the argument that consumes it stays inside the counted proof span.\par\medskip
\begin{lemma}
\label{lemma-restriction-derived-complete-equivalence}
Let $A \to B$ be a ring map. Let $I \subset A$ be a finitely generated ideal.
If $A \to B$ is flat and $A/I \cong B/IB$, then the restriction functor
$D(B) \to D(A)$ induces an equivalence
$D_{comp}(B, IB) \to D_{comp}(A, I)$.
\end{lemma}

\begin{proof}
Choose generators $f_1, \ldots, f_r$ of $I$.
Denote $\check{\mathcal{C}}^\bullet_A \to \check{\mathcal{C}}^\bullet_B$
the quasi-isomorphism of extended alternating {\v C}ech complexes of
Lemma \ref{lemma-map-identifies-koszul-and-cech-complexes}.
Let $K \in D_{comp}(A, I)$. Let $I^\bullet$ be a K-injective
complex of $A$-modules representing $K$. Since $\Ext^n_A(A_f, K)$
and $\Ext^n_A(B_f, K)$ are zero for all $f \in I$ and
$n \in \mathbf{Z}$ (Lemma \ref{lemma-hom-from-Af}) we conclude that
$\check{\mathcal{C}}^\bullet_A \to A$ and
$\check{\mathcal{C}}^\bullet_B \to B$ induce quasi-isomorphisms
$$
I^\bullet = \Hom_A(A, I^\bullet) \longrightarrow
\text{Tot}(\Hom_A(\check{\mathcal{C}}^\bullet_A, I^\bullet))
$$
and
$$
\Hom_A(B, I^\bullet) \longrightarrow
\text{Tot}(\Hom_A(\check{\mathcal{C}}^\bullet_B, I^\bullet))
$$
Some details omitted.
Since $\check{\mathcal{C}}^\bullet_A \to \check{\mathcal{C}}^\bullet_B$
is a quasi-isomorphism and $I^\bullet$ is K-injective we conclude
that $\Hom_A(B, I^\bullet) \to I^\bullet$ is a quasi-isomorphism.
As the complex $\Hom_A(B, I^\bullet)$ is a complex of $B$-modules
we conclude that $K$ is in the image of the restriction map, i.e.,
the functor is essentially surjective

\medskip\noindent
In fact, the argument shows that
$F : D_{comp}(A, I) \to D_{comp}(B, IB)$, $K \mapsto \Hom_A(B, I^\bullet)$
is a left inverse to restriction. Finally, suppose that
$L \in D_{comp}(B, IB)$. Represent $L$ by a K-injective complex
$J^\bullet$ of $B$-modules.
Then $J^\bullet$ is also K-injective as a complex of $A$-modules
(Lemma \ref{lemma-K-injective-flat}) hence
$F(\text{restriction of }L) = \Hom_A(B, J^\bullet)$.
There is a map $J^\bullet \to \Hom_A(B, J^\bullet)$
of complexes of $B$-modules, whose composition with
$\Hom_A(B, J^\bullet) \to J^\bullet$ is the identity.
We conclude that $F$ is also a right inverse to restriction
and the proof is finished.
\end{proof}

\end{document}
